\sectionheading{5.1 প্রোগ্রামের ধারণা}

কম্পিউটার অত্যন্ত দ্রুত ও নির্ভুলভাবে কাজ করতে পারলেও
\textbf{কী কাজ করতে হবে তা নিজে থেকে নির্ধারণ করতে পারে না}।
কোনো কাজ সম্পন্ন করতে তাকে সঠিক ক্রমে প্রয়োজনীয় নির্দেশ দিতে হয়।
এই নির্দেশনাগুলো পরিকল্পিত ও সুশৃঙ্খলভাবে সাজিয়ে তৈরি করা হয়
\textbf{Program}।

সহজ কথায়, কম্পিউটারকে যদি কোনো Problem দেওয়া হয়, তবে সেই Problem
সমাধানের জন্য কী কী কাজ কোন ক্রমে করতে হবে—Program-এ সেটিই
নির্দেশ আকারে লেখা থাকে।

%================================================
% Definition
%================================================
\booksubheading{প্রোগ্রাম কাকে বলে?}

\begin{definitionbox}

\textbf{\color{darkTeal}সংজ্ঞা}

কোনো নির্দিষ্ট সমস্যা সমাধান বা কাজ সম্পাদনের জন্য কম্পিউটারকে
দেওয়া \textbf{যুক্তিসংগত, সুশৃঙ্খল ও ধারাবাহিক নির্দেশনার সমষ্টিকে}
\textbf{প্রোগ্রাম (Program)} বলে।

\end{definitionbox}

\textbf{সহজভাবে মনে রাখবে—}

\begin{center}
\textbf{
সমস্যা সমাধানের জন্য সাজানো নির্দেশনা
\(\;=\;\)
Program
}
\end{center}

%================================================
% Program works
%================================================
\booksubheading{একটি Program কীভাবে কাজ করে?}

একটি Program-এর কাজ সাধারণত তিনটি মৌলিক পর্যায়ে বোঝানো যায়—
\textbf{Input, Processing ও Output}।

\begin{center}

\begin{tikzpicture}[node distance=12mm]

\node[databox] (input)
{\textbf{Input}\\Data গ্রহণ};

\node[processbox,right=of input] (process)
{\textbf{Processing}\\নির্দেশ অনুযায়ী কাজ};

\node[databox,right=of process] (output)
{\textbf{Output}\\ফলাফল প্রদান};

\draw[arrow] (input) -- (process);
\draw[arrow] (process) -- (output);

\end{tikzpicture}

\end{center}

\begin{itemize}

\item \textbf{Input:} Program-এর কাজ করার জন্য প্রয়োজনীয় Data গ্রহণ করা।

\item \textbf{Processing:} নির্ধারিত Formula, Rule বা Logic অনুসারে
Input Data-এর ওপর কাজ করা।

\item \textbf{Output:} Processing শেষে পাওয়া Result ব্যবহারকারীর
কাছে প্রদর্শন করা।

\end{itemize}

%================================================
% Example
%================================================
\begin{examplebox}{উদাহরণ: দুটি সংখ্যার যোগফল}

ধরা যাক, দুটি সংখ্যা \textbf{10} ও \textbf{20}-এর যোগফল বের করতে হবে।

এক্ষেত্রে Program-এর কাজ হবে—

\begin{itemize}

\item 10 ও 20 সংখ্যা দুটি গ্রহণ করা।

\item সংখ্যা দুটি যোগ করা।

\item পাওয়া যোগফল প্রদর্শন করা।

\end{itemize}

\[
\boxed{
10,\;20
\;\longrightarrow\;
10+20
\;\longrightarrow\;
30
}
\]

অর্থাৎ—

\begin{center}
\begin{tabularx}{0.78\linewidth}{|>{\bfseries}C{0.22\linewidth}|Y|}
\hline
Input & 10, 20\\
\hline
Processing & \(10+20\)\\
\hline
Output & 30\\
\hline
\end{tabularx}
\end{center}

\end{examplebox}

%================================================
% Program and problem
%================================================
\booksubheading{Problem থেকে Program}

Program লেখার শুরুতেই Code লেখা হয় না। প্রথমে Problem বুঝতে হয়,
তারপর সমাধানের ধাপ নির্ধারণ করে সেই ধাপগুলো Programming Language-এর
সাহায্যে Program আকারে লেখা হয়।

\begin{center}

\textbf{
Problem
\(\rightarrow\)
সমাধানের চিন্তা
\(\rightarrow\)
নির্দেশনা
\(\rightarrow\)
Program
\(\rightarrow\)
Result
}

\end{center}

যেমন, একজন শিক্ষার্থীর তিনটি বিষয়ের গড় নম্বর বের করতে হলে প্রথমে
তিনটি নম্বর Input নিতে হবে, নম্বরগুলো যোগ করতে হবে, যোগফলকে 3 দিয়ে
ভাগ করতে হবে এবং শেষে Average প্রদর্শন করতে হবে।

%================================================
% Second example
%================================================
\begin{examplebox}{আরও একটি উদাহরণ: গড় নম্বর নির্ণয়}

তিনটি বিষয়ের নম্বর—

\[
70,\qquad80,\qquad90
\]

প্রথমে মোট নম্বর—

\[
Total=70+80+90=240
\]

তারপর গড়—

\[
Average=\frac{240}{3}=80
\]

অতএব,

\[
\boxed{Average=80}
\]

এখানে \textbf{70, 80, 90 হলো Input}, যোগ ও ভাগ হলো
\textbf{Processing} এবং \textbf{80 হলো Output}।

\end{examplebox}

%================================================
% Fundamental parts
%================================================
\booksubheading{Program-এর মৌলিক অংশ}

\begin{tabularx}{\linewidth}{|>{\bfseries}p{0.18\linewidth}|Y|Y|}
\hline

\rowcolor{softBlue}
\multicolumn{1}{|c|}{\textbf{অংশ}}
&
\multicolumn{1}{c|}{\textbf{কাজ}}
&
\multicolumn{1}{c|}{\textbf{উদাহরণ}}
\\
\hline

Input
&
প্রয়োজনীয় Data গ্রহণ করা
&
দুটি সংখ্যা গ্রহণ
\\
\hline

Processing
&
Data-এর ওপর প্রয়োজনীয় Logic বা Operation প্রয়োগ করা
&
সংখ্যা দুটি যোগ করা
\\
\hline

Output
&
Process শেষে পাওয়া Result প্রদর্শন করা
&
যোগফল প্রদর্শন
\\
\hline

\end{tabularx}

%================================================
% Instructions
%================================================
\booksubheading{Program-এর নির্দেশনা কেমন হওয়া উচিত?}

Program-এর Instruction এলোমেলো হলে কম্পিউটার কাঙ্ক্ষিত Result
দিতে পারবে না। তাই নির্দেশগুলো কিছু সাধারণ বৈশিষ্ট্য অনুসরণ করা
প্রয়োজন।

\begin{itemize}

\item \textbf{সুনির্দিষ্ট হতে হবে—} কোন কাজ করতে হবে তা পরিষ্কার থাকতে হবে।

\item \textbf{সঠিক ক্রমে থাকতে হবে—} যে কাজ আগে করা দরকার সেটি আগে
এবং পরের কাজ পরে থাকতে হবে।

\item \textbf{যুক্তিসংগত হতে হবে—} প্রতিটি Instruction Problem সমাধানের
সঙ্গে সম্পর্কিত হতে হবে।

\item \textbf{সম্পূর্ণ হতে হবে—} প্রয়োজনীয় কোনো ধাপ বাদ দেওয়া যাবে না।

\item \textbf{কার্যকর হতে হবে—} নির্দেশগুলো এমন হতে হবে যাতে বাস্তবে
কম্পিউটার তা সম্পন্ন করতে পারে।

\end{itemize}

\begin{examplebox}{নির্দেশের ক্রম কেন গুরুত্বপূর্ণ?}

চা তৈরির ক্ষেত্রে যেমন—

\begin{center}
পানি গরম
\(\rightarrow\)
চা-পাতা
\(\rightarrow\)
চিনি/দুধ
\(\rightarrow\)
চা প্রস্তুত
\end{center}

ধাপগুলোর ক্রম এলোমেলো হলে কাঙ্ক্ষিত ফল নাও পাওয়া যেতে পারে।
Program-এর ক্ষেত্রেও Instruction-এর সঠিক ক্রম অত্যন্ত গুরুত্বপূর্ণ।

\end{examplebox}

%================================================
% Importance
%================================================
\booksubheading{Program কেন প্রয়োজন?}

কম্পিউটারের Hardware শক্তিশালী হলেও Program ছাড়া তা নির্দিষ্ট
ব্যবহারকারীর কাজ সম্পন্ন করতে পারে না। Program-এর সাহায্যে—

\begin{compacttwocolpoints}

\item গাণিতিক ও যৌক্তিক হিসাব দ্রুত করা যায়।

\item একই কাজ বারবার স্বয়ংক্রিয়ভাবে সম্পন্ন করা যায়।

\item বিপুল পরিমাণ Data দ্রুত Process করা যায়।

\item শিক্ষাপ্রতিষ্ঠানের Result, Bank হিসাব, Billing System,
Reservation System ইত্যাদি পরিচালনা করা যায়।

\item Website, Mobile App, Game ও বিভিন্ন Software তৈরি করা যায়।

\item সময়, শ্রম ও অনেক ক্ষেত্রে ভুলের পরিমাণ কমানো যায়।

\end{compacttwocolpoints}

%================================================
% Good Program
%================================================
\booksubheading{একটি ভালো Program-এর সাধারণ বৈশিষ্ট্য}

\begin{tabularx}{\linewidth}{|>{\bfseries}p{0.23\linewidth}|Y|}
\hline

\rowcolor{softBlue}
\multicolumn{1}{|c|}{\textbf{বৈশিষ্ট্য}}
&
\multicolumn{1}{c|}{\textbf{ধারণা}}
\\
\hline

Correctness
&
সঠিক Input-এর জন্য সঠিক Result দিতে হবে।
\\
\hline

Clarity
&
Program-এর Logic ও Instruction সহজে বোঝা যায় এমন হওয়া উচিত।
\\
\hline

Efficiency
&
অপ্রয়োজনীয় ধাপ ও Resource-এর ব্যবহার যথাসম্ভব কম হওয়া উচিত।
\\
\hline

Reliability
&
বিভিন্ন বৈধ Input-এর ক্ষেত্রেও নির্ভরযোগ্যভাবে কাজ করা উচিত।
\\
\hline

Maintainability
&
প্রয়োজনে Program সহজে সংশোধন ও উন্নত করা সম্ভব হওয়া উচিত।
\\
\hline

\end{tabularx}

%================================================
% Program and Programming Language
%================================================
\booksubheading{Program ও Programming Language-এর সম্পর্ক}

Program নিজে কোনো ভাষা নয়। Program হলো নির্দিষ্ট কাজের জন্য
লেখা Instruction-এর সমষ্টি, আর সেই Instruction লেখার জন্য
Programming Language ব্যবহার করা হয়।

\begin{center}

\textbf{
Programming Language
\(\longrightarrow\)
Program লেখা
\(\longrightarrow\)
কম্পিউটারকে কাজের নির্দেশ
}

\end{center}

\begin{rememberbox}

\textbf{\color{deepOrange}সহজভাবে মনে রাখবে}

\textbf{Programming Language হলো Program লেখার মাধ্যম, আর Program
হলো সেই ভাষায় লেখা কার্যনির্দেশ।}

\end{rememberbox}

%================================================
% Program vs Software
%================================================
\booksubheading{Program ও Software কি একই?}

Program ও Software পরস্পর সম্পর্কিত হলেও সবসময় একই বিষয় নয়।

\begin{tabularx}{\linewidth}{|Y|Y|}
\hline

\rowcolor{softBlue}
\multicolumn{1}{|c|}{\textbf{Program}}
&
\multicolumn{1}{c|}{\textbf{Software}}
\\
\hline

সাধারণত নির্দিষ্ট কাজ সম্পাদনের Instruction-এর সমষ্টি
&
এক বা একাধিক Program-এর সঙ্গে প্রয়োজনীয় Data, Library,
Interface ইত্যাদির সমন্বয়ে পূর্ণাঙ্গ ব্যবস্থা হতে পারে
\\
\hline

একটি ছোট বা নির্দিষ্ট কাজ করতে পারে
&
একাধিক সম্পর্কিত কাজ সম্পন্ন করতে পারে
\\
\hline

উদাহরণ: দুটি সংখ্যার যোগফল নির্ণয়ের Program
&
উদাহরণ: Word Processing Software
\\
\hline

\end{tabularx}

%================================================
% Final reminder
%================================================
\begin{rememberbox}

\textbf{\color{deepOrange}মনে রাখবে}

\begin{itemize}

\item Program হলো \textbf{নির্দেশনার সমষ্টি}।

\item Program-এর সাধারণ কাজের ধারা—
\textbf{Input \(\rightarrow\) Processing \(\rightarrow\) Output}।

\item Instruction-এর \textbf{সঠিক ক্রম} Program-এর Result-এর জন্য গুরুত্বপূর্ণ।

\item কোনো Problem দেওয়া হলে তার \textbf{Input, Processing ও Output}
আলাদা করে শনাক্ত করতে পারা উচিত।

\item Program লেখার জন্য Programming Language প্রয়োজন।

\end{itemize}

\end{rememberbox}

\small

%================================================
% 5.2 Programming Language
%================================================
